\documentclass{ximera}

\input{../../../preamble.tex}


\definecolor{fvdnavy}{HTML}{071D43}
\definecolor{fvdcyan}{HTML}{20BFC6}
\definecolor{fvdcyanlight}{HTML}{EFFCFB}
\definecolor{fvdmagenta}{HTML}{F10D69}
\definecolor{fvdmagentalight}{HTML}{FFF1F7}
\definecolor{fvdtext}{HTML}{163044}





%=================================================
% Pedagogische omgevingen
% PDF: tcolorbox
% HTML: learning-box
%=================================================

\newenvironment{voorkennis}
{%
    \ifdefined\HCode
        \HCode{%
<section class="learning-box learning-box--prior">
<div class="learning-box__header">
<span class="learning-box__icon" aria-hidden="true"></span>
<span class="learning-box__title">Voorkennis</span>
</div>
<div class="learning-box__content">
        }%
        \begin{itemize}
    \else
       \begin{tcolorbox}[
    enhanced,
    breakable,
    colback=fvdcyanlight,
    colframe=fvdcyan,
    coltext=fvdtext,
    boxrule=0.5pt,
    leftrule=4pt,
    arc=4mm,
    outer arc=4mm,
    left=7mm,
    right=6mm,
    top=5mm,
    bottom=5mm,
    before skip=6mm,
    after skip=6mm,
    title=Voorkennis,
    coltitle=fvdcyan,
    fonttitle=\bfseries\Large,
    attach title to upper={\par\smallskip},
    boxed title style={
        empty,
        left=0mm,
        right=0mm,
        top=0mm,
        bottom=0mm
    },
    overlay={
        \draw[
            fvdcyan,
            line width=0.5pt
        ]
        ([xshift=42mm,yshift=-13mm]frame.north west)
        --
        ([xshift=-7mm,yshift=-13mm]frame.north east);
    }
]
        \begin{itemize}
    \fi
}
{%
    \end{itemize}
    \ifdefined\HCode
        \HCode{%
</div>
</section>
        }%
    \else
        \end{tcolorbox}
    \fi
}


\newenvironment{leerdoelen}
{%
    \ifdefined\HCode
        \HCode{%
<section class="learning-box learning-box--goals">
<div class="learning-box__header">
<span class="learning-box__icon" aria-hidden="true"></span>
<span class="learning-box__title">Leerdoelen</span>
</div>
<div class="learning-box__content">
        }%
        \begin{itemize}
    \else
\begin{tcolorbox}[
    enhanced,
    breakable,
    colback=fvdmagentalight,
    colframe=fvdmagenta,
    coltext=fvdtext,
    boxrule=0.5pt,
    leftrule=4pt,
    arc=4mm,
    outer arc=4mm,
    left=7mm,
    right=6mm,
    top=5mm,
    bottom=5mm,
    before skip=6mm,
    after skip=6mm,
    title=Leerdoelen,
    coltitle=fvdmagenta,
    fonttitle=\bfseries\Large,
    attach title to upper={\par\smallskip},
    boxed title style={
        empty,
        left=0mm,
        right=0mm,
        top=0mm,
        bottom=0mm
    },
    overlay={
        \draw[
            fvdmagenta,
            line width=0.5pt
        ]
        ([xshift=42mm,yshift=-13mm]frame.north west)
        --
        ([xshift=-7mm,yshift=-13mm]frame.north east);
    }
]
        \begin{itemize}
    \fi
}
{%
    \end{itemize}
    \ifdefined\HCode
        \HCode{%
</div>
</section>
        }%
    \else
        \end{tcolorbox}
    \fi
}


\addPrintStyle{../../..}

\title{Wiskundige uitspraken}
\author{Ingmar Herreman}

\begin{abstract}
Wiskunde bestaat niet alleen uit berekeningen en formules.
We doen ook voortdurend \textbf{uitspraken}: beweringen waarvan we
willen weten of ze waar of onwaar zijn.

In dit hoofdstuk onderzoeken we wat een wiskundige uitspraak precies is.
We leren uitspraken onderscheiden van vragen, opdrachten en formules
met veranderlijken. Daarbij ontdekken we dat de betekenis van een
wiskundige bewering ook kan afhangen van het domein waarin we werken.

Vervolgens maken we kennis met \textbf{open uitspraken} en
\textbf{predicaten}. Daarmee kunnen we eigenschappen van getallen
en andere wiskundige objecten nauwkeurig formuleren.

Ten slotte onderzoeken we waarom voorbeelden nuttig zijn, maar niet
volstaan om een algemene uitspraak te bewijzen, en hoe één goed gekozen
tegenvoorbeeld een algemene bewering wel kan weerleggen.

Zo leggen we het fundament voor de logische taal waarmee we in de
volgende hoofdstukken samengestelde uitspraken en wiskundige
redeneringen zullen analyseren.
\end{abstract}

\begin{document}

% FVD_CHAPTER_DATA_START
% Automatisch gegenereerd uit index.tex.
% Niet handmatig aanpassen.
\fvdchapterdata{1}{Logica — De taal van correct redeneren}{1}
% FVD_CHAPTER_DATA_END













\maketitle


% ============================================================
% VOORKENNIS
% ============================================================

\begin{voorkennis}
  \item Je kan werken met de getallenverzamelingen
        $\mathbb{N}$, $\mathbb{Z}$, $\mathbb{Q}$ en $\mathbb{R}$.
  \item Je kan eenvoudige gelijkheden en ongelijkheden interpreteren.
  \item Je kan waarden invullen voor een veranderlijke.
  \item Je kent begrippen zoals even en oneven getal, veelvoud en priemgetal.
  \item Je kan een wiskundige bewering controleren aan de hand van een berekening of voorbeeld.
\end{voorkennis}


% ============================================================
% LEERDOELEN
% ============================================================

\begin{leerdoelen}
  \item Je kan herkennen of een zin een wiskundige uitspraak is.
  \item Je kan de waarheidswaarde van een eenvoudige wiskundige uitspraak bepalen en verantwoorden.
  \item Je kan uitleggen waarom vragen en opdrachten geen wiskundige uitspraken zijn.
  \item Je kan een uitspraak onderscheiden van een open uitspraak.
  \item Je kan een predicaat herkennen en met een notatie zoals $P(x)$ beschrijven.
  \item Je kan onderzoeken voor welke waarden een predicaat waar of onwaar is.
  \item Je kan uitleggen waarom het gekozen domein belangrijk is voor de betekenis van een uitspraak.
  \item Je kan een open uitspraak omvormen tot een uitspraak door een veranderlijke een waarde te geven.
  \item Je kan een algemene bewering onderzoeken met voorbeelden en tegenvoorbeelden.
  \item Je kan uitleggen waarom voorbeelden een algemene uitspraak niet bewijzen en waarom één tegenvoorbeeld ze wel kan weerleggen.
\end{leerdoelen}


% ============================================================
\FVDsection{Wanneer is iets een wiskundige uitspraak?}
% ============================================================

Bekijk de volgende zinnen.

\[
\begin{array}{ll}
\text{A.} & 17\text{ is een priemgetal.}\\[2mm]
\text{B.} & 15\text{ is een priemgetal.}\\[2mm]
\text{C.} & \text{Is }17\text{ een priemgetal?}\\[2mm]
\text{D.} & \text{Bereken }17+25.\\[2mm]
\text{E.} & x+3=8.
\end{array}
\]

Ze hebben allemaal met wiskunde te maken, maar logisch gezien zijn ze
niet van hetzelfde type.

Van zin A kunnen we bepalen dat hij \textbf{waar} is.

Zin B is \textbf{onwaar}.

Zin C is een vraag en zin D is een opdracht. Ze beweren zelf niets
waarvan we kunnen zeggen dat het waar of onwaar is.

Bij zin E ontstaat nog een ander probleem. Voor $x=5$ is de gelijkheid
waar, maar voor $x=2$ is ze onwaar. Zonder te weten wat $x$ voorstelt,
heeft de formule dus nog geen vaste waarheidswaarde.

Dit onderscheid vormt het vertrekpunt van de wiskundige logica.


\begin{definitie}
Een \textbf{wiskundige uitspraak} is een bewering waarvan
ondubbelzinnig kan worden vastgesteld of ze \textbf{waar} of
\textbf{onwaar} is.

De waarde \emph{waar} of \emph{onwaar} noemen we de
\textbf{waarheidswaarde} van de uitspraak.
\end{definitie}


\begin{voorbeeld}[Ware en onware uitspraken]

De bewering

\[
8+5=13
\]

is een wiskundige uitspraak en is waar.

Ook

\[
7^2=42
\]

is een wiskundige uitspraak, maar ze is onwaar, want

\[
7^2=49.
\]

Het feit dat een bewering onwaar is, verhindert dus niet dat ze een
wiskundige uitspraak is.
\end{voorbeeld}


\begin{opmerking}
Een \textbf{uitspraak} is niet hetzelfde als een
\textbf{ware uitspraak}.

Zowel

\[
2+3=5
\]

als

\[
2+3=7
\]

zijn uitspraken. Ze hebben alleen een verschillende waarheidswaarde.
\end{opmerking}


% ------------------------------------------------------------
\FVDsubsection{Niet elke wiskundige zin is een uitspraak}
% ------------------------------------------------------------

Een zin kan wiskundige symbolen of begrippen bevatten zonder zelf een
wiskundige uitspraak te zijn.

Vergelijk:

\[
\begin{array}{ll}
\text{A.} & 25\text{ is een kwadraatgetal.}\\
\text{B.} & \text{Is }25\text{ een kwadraatgetal?}\\
\text{C.} & \text{Bereken }\sqrt{25}.
\end{array}
\]

Alleen A is een uitspraak.

B is een \textbf{vraag} en C een \textbf{opdracht}.


\begin{oefening}[Uitspraak of niet?]{\oefB \ESS}

Bepaal telkens of de zin een wiskundige uitspraak is.

\begin{question}
\[
11\text{ is een priemgetal.}
\]

\begin{multipleChoice}
  \choice[correct]{Dit is een uitspraak.}
  \choice{Dit is geen uitspraak.}
\end{multipleChoice}
\end{question}

\begin{question}
\[
\text{Bereken }3^4.
\]

\begin{multipleChoice}
  \choice{Dit is een uitspraak.}
  \choice[correct]{Dit is geen uitspraak.}
\end{multipleChoice}
\end{question}

\begin{question}
\[
\sqrt{2}\in\mathbb{Q}.
\]

\begin{multipleChoice}
  \choice[correct]{Dit is een uitspraak.}
  \choice{Dit is geen uitspraak.}
\end{multipleChoice}
\end{question}

\begin{question}
Leg uit waarom de vorige bewering een uitspraak is, ook al is ze onwaar.

\begin{oplossing}
De bewering heeft een ondubbelzinnige waarheidswaarde.
Omdat $\sqrt{2}$ niet rationaal is, is de uitspraak onwaar.
Een uitspraak hoeft niet waar te zijn om een uitspraak te zijn.
\end{oplossing}
\end{question}

\end{oefening}


% ============================================================
\FVDsection{Wiskundige taal moet precies zijn}
% ============================================================

In gewone taal gebruiken we vaak woorden waarvan de betekenis afhangt
van de context.

Beschouw bijvoorbeeld:

\begin{center}
\textit{20 is een groot getal.}
\end{center}

Is dat waar?

Dat hangt af van wat we met \emph{groot} bedoelen.

Voor het aantal leerlingen in een klas kan $20$ een behoorlijk aantal
zijn. Voor het aantal moleculen in een druppel water is $20$ juist
uitzonderlijk klein.

Vergelijk dit met

\[
20>15.
\]

Hier is geen interpretatie nodig. De bewering heeft een
ondubbelzinnige waarheidswaarde.


\begin{opmerking}
Wiskundige logica vereist dat we zo precies mogelijk formuleren
\textbf{wat} we beweren en \textbf{binnen welke context} we dat
beweren.

Dat wordt later essentieel wanneer we willen bepalen of een conclusie
werkelijk uit bepaalde voorwaarden volgt.
\end{opmerking}


\begin{voorbeeld}[Van vaag naar precies]

De zin

\begin{center}
\textit{$x$ is dicht bij $10$.}
\end{center}

is zonder verdere afspraak niet precies genoeg.

We kunnen de bedoeling bijvoorbeeld nauwkeuriger formuleren als

\[
|x-10|<0{,}1.
\]

Nu ligt exact vast welke waarden van $x$ aan de voorwaarde voldoen.
\end{voorbeeld}


% ============================================================
\FVDsection{Uitspraken benoemen}
% ============================================================

Wanneer we later verschillende uitspraken met elkaar verbinden,
zou het omslachtig zijn om telkens de volledige zin te schrijven.

Daarom geven we uitspraken vaak een korte naam.

We gebruiken daarvoor meestal letters zoals

\[
p,\qquad q,\qquad r,\ldots
\]


\begin{voorbeeld}[Twee uitspraken]

We definiëren

\[
p:\quad 12\text{ is even}
\]

en

\[
q:\quad 12\text{ is deelbaar door }3.
\]

Beide uitspraken zijn waar.

De letters $p$ en $q$ stellen hier dus geen getallen voor.
Elke letter verwijst naar een \textbf{volledige uitspraak}.
\end{voorbeeld}


\begin{opmerking}
In de volgende hoofdstukken zullen we uitspraken combineren.

Zo zullen we bijvoorbeeld onderzoeken wat precies bedoeld wordt met

\[
p\text{ en }q,
\qquad
p\text{ of }q
\]

en

\[
\text{als }p\text{, dan }q.
\]

Daarvoor moeten we eerst goed begrijpen wat $p$ en $q$ afzonderlijk
betekenen.
\end{opmerking}


% ============================================================
\FVDsection{Open uitspraken}
% ============================================================

Beschouw nu

\[
x+3=8.
\]

Kunnen we zeggen dat deze formule waar of onwaar is?

Voor

\[
x=5
\]

krijgen we

\[
5+3=8,
\]

en is de bewering waar.

Voor

\[
x=2
\]

krijgen we

\[
2+3=8,
\]

en is de bewering onwaar.

De waarheidswaarde hangt dus af van de gekozen waarde van $x$.


\begin{definitie}
Een \textbf{open uitspraak} is een bewering met één of meer
veranderlijken waarvan de waarheidswaarde afhangt van de waarden
van die veranderlijken.
\end{definitie}


\begin{voorbeeld}[Een open uitspraak]

Beschouw

\[
x^2=9.
\]

Voor $x=3$ is de bewering waar.

Voor $x=-3$ is ze eveneens waar.

Voor $x=2$ is ze onwaar.

Zolang $x$ geen bepaalde waarde heeft, is $x^2=9$ dus een open
uitspraak.
\end{voorbeeld}


\begin{oefening}[Van open uitspraak naar uitspraak]{\oefB \ESS}

Beschouw

\[
x^2<10.
\]

\begin{question}
Is de formule zonder verdere informatie een uitspraak met één vaste
waarheidswaarde?

\begin{multipleChoice}
  \choice{Ja.}
  \choice[correct]{Nee.}
\end{multipleChoice}
\end{question}

\begin{question}
Onderzoek de formule voor $x=3$.

\begin{oplossing}
We krijgen

\[
3^2<10
\]

of

\[
9<10.
\]

De ontstane uitspraak is waar.
\end{oplossing}
\end{question}

\begin{question}
Onderzoek de formule voor $x=4$.

\begin{oplossing}
We krijgen

\[
4^2<10
\]

of

\[
16<10.
\]

De ontstane uitspraak is onwaar.
\end{oplossing}
\end{question}

\end{oefening}


% ============================================================
\FVDsection{Predicaten}
% ============================================================

Open uitspraken komen in de wiskunde voortdurend voor.

Om er gemakkelijk naar te verwijzen, gebruiken we een notatie zoals

\[
P(x).
\]

Beschouw bijvoorbeeld

\[
P(x):\quad x^2>4.
\]

De letter $P$ staat voor de eigenschap

\begin{center}
\textit{het kwadraat van $x$ is groter dan $4$.}
\end{center}


\begin{definitie}
Een \textbf{predicaat} $P(x)$ beschrijft een eigenschap van $x$.

Wanneer we voor $x$ een toegelaten waarde invullen, ontstaat een
uitspraak waarvan we de waarheidswaarde kunnen bepalen.
\end{definitie}


\begin{voorbeeld}[Even zijn als predicaat]

Definieer

\[
P(n):\quad n\text{ is even}.
\]

Dan is

\[
P(8)
\]

waar, terwijl

\[
P(11)
\]

onwaar is.
\end{voorbeeld}


Een predicaat kan ook meer dan één veranderlijke bevatten.


\begin{voorbeeld}[Een predicaat met twee veranderlijken]

Definieer

\[
Q(x,y):\quad x<y.
\]

Dan is

\[
Q(2,5)
\]

waar, want $2<5$.

Daarentegen is

\[
Q(7,3)
\]

onwaar.
\end{voorbeeld}


\begin{opmerking}
Let op het verschil tussen

\[
p:\quad 7\text{ is een priemgetal}
\]

en

\[
P(x):\quad x\text{ is een priemgetal}.
\]

De eerste formule stelt één concrete uitspraak voor.

De tweede beschrijft een eigenschap die we voor verschillende waarden
van $x$ kunnen onderzoeken.
\end{opmerking}


% ============================================================
\FVDsection{Het domein doet ertoe}
% ============================================================

Wanneer een veranderlijke voorkomt, moeten we weten welke waarden die
veranderlijke mag aannemen.

Beschouw bijvoorbeeld

\[
P(x):\quad x^2=2.
\]

Als we alleen natuurlijke getallen toelaten, dan is er geen enkele
waarde van $x$ waarvoor $P(x)$ waar is.

Werken we in de reële getallen, dan zijn er wel oplossingen:

\[
x=\sqrt{2}
\qquad\text{en}\qquad
x=-\sqrt{2}.
\]

De verzameling waarin we werken speelt dus een essentiële rol.


\begin{definitie}
Het \textbf{domein} of \textbf{universum} van een veranderlijke is
de verzameling van waarden die de veranderlijke mag aannemen.
\end{definitie}


\begin{voorbeeld}[Hetzelfde predicaat, een ander domein]

Beschouw

\[
P(x):\quad x^2=4.
\]

In het domein

\[
\mathbb{N}
\]

is $x=2$ een waarde waarvoor $P(x)$ waar is.

In het domein

\[
\mathbb{Z}
\]

zijn zowel $x=2$ als $x=-2$ waarden waarvoor $P(x)$ waar is.

Het predicaat is hetzelfde, maar de toegelaten waarden hangen af
van het domein.
\end{voorbeeld}


\begin{voorbeeld}[Een vergelijking zonder oplossing in elk domein]

Beschouw

\[
Q(x):\quad x^2=-1.
\]

Voor

\[
x\in\mathbb{R}
\]

is er geen waarde waarvoor $Q(x)$ waar is.

Wanneer we later met complexe getallen werken, verandert die situatie.

Ook dit toont waarom het domein deel uitmaakt van de betekenis van
een wiskundige bewering.
\end{voorbeeld}


\begin{opmerking}
Wanneer je een uitspraak met veranderlijken analyseert, stel dan steeds
de volgende vragen:

\[
\important{
\begin{array}{c}
\text{Welke veranderlijke komt voor?}\\
\text{Welke waarden mag ze aannemen?}\\
\text{In welk domein werken we?}
\end{array}
}
\]
\end{opmerking}


% ============================================================
\FVDsection{Van een open uitspraak naar een uitspraak}
% ============================================================

Een open uitspraak heeft nog geen vaste waarheidswaarde.

We hebben al gezien dat we dit kunnen veranderen door een concrete
waarde voor de veranderlijke in te vullen.


\begin{voorbeeld}

Gegeven is

\[
P(x):\quad x^2<25,
\qquad x\in\mathbb{R}.
\]

Voor $x=3$ krijgen we

\[
P(3):\quad 3^2<25,
\]

dus een ware uitspraak.

Voor $x=7$ krijgen we

\[
P(7):\quad 7^2<25,
\]

dus een onware uitspraak.
\end{voorbeeld}


Maar in de wiskunde willen we vaak veel meer zeggen.

We willen bijvoorbeeld beweren:

\begin{quote}
Voor elk reëel getal geldt dat zijn kwadraat niet-negatief is.
\end{quote}

Of:

\begin{quote}
Er bestaat een reëel getal waarvan het kwadraat gelijk is aan $2$.
\end{quote}

De woorden

\begin{center}
\textbf{voor elk}
\qquad\text{en}\qquad
\textbf{er bestaat}
\end{center}

maken van een open uitspraak opnieuw een uitspraak met een
waarheidswaarde.

Later zullen we hiervoor de symbolen

\[
\forall
\qquad\text{en}\qquad
\exists
\]

gebruiken.

Deze symbolen noemen we \textbf{kwantoren}.


\begin{opmerking}
We werken kwantoren hier nog niet systematisch uit.

Voorlopig is vooral het volgende onderscheid belangrijk:

\[
\important{
\begin{array}{c}
P(x) \text{ beschrijft een eigenschap van }x,\\[1mm]
P(a) \text{ is na het invullen van }a\text{ een concrete uitspraak.}
\end{array}
}
\]
\end{opmerking}


% ============================================================
\FVDsection{Voorbeelden onderzoeken}
% ============================================================

Wanneer we een wiskundige bewering tegenkomen, is het vaak verstandig
om eerst enkele voorbeelden te onderzoeken.

Beschouw bijvoorbeeld de bewering:

\begin{center}
\textit{Het kwadraat van een oneven geheel getal is oneven.}
\end{center}

We kunnen controleren:

\[
3^2=9,
\qquad
5^2=25,
\qquad
7^2=49,
\qquad
11^2=121.
\]

Al deze voorbeelden ondersteunen de bewering.

Maar hebben we daarmee bewezen dat ze voor \textbf{elk} oneven geheel
getal geldt?

Nee.

We hebben slechts een eindig aantal gevallen gecontroleerd.


\begin{opmerking}
Voorbeelden zijn bijzonder nuttig om

\begin{itemize}
  \item een uitspraak te begrijpen;
  \item patronen te ontdekken;
  \item een vermoeden te formuleren;
  \item mogelijke fouten op te sporen.
\end{itemize}

Maar een verzameling voorbeelden vormt in het algemeen geen bewijs
van een uitspraak die over alle mogelijke gevallen gaat.
\end{opmerking}


Dit onderscheid zal later belangrijk worden wanneer we wiskundige
bewijzen leren opbouwen.


% ============================================================
\FVDsection{Tegenvoorbeelden}
% ============================================================

Bekijk nu de volgende bewering:

\begin{center}
\textit{Alle priemgetallen zijn oneven.}
\end{center}

We controleren enkele priemgetallen:

\[
3,\quad5,\quad7,\quad11,\quad13,\quad17.
\]

Ze zijn allemaal oneven.

Toch is de bewering onwaar.

Het getal

\[
2
\]

is namelijk zowel priem als even.


\begin{definitie}
Een \textbf{tegenvoorbeeld} is een voorbeeld dat aantoont dat een
algemene bewering niet waar is.
\end{definitie}


In dit geval is

\[
2
\]

een tegenvoorbeeld voor de bewering

\begin{center}
\textit{Alle priemgetallen zijn oneven.}
\end{center}


Daar ontstaat een belangrijke asymmetrie:

\[
\important{
\begin{array}{c}
\text{Veel bevestigende voorbeelden bewijzen een algemene uitspraak niet.}\\[2mm]
\text{Eén geldig tegenvoorbeeld volstaat om ze te weerleggen.}
\end{array}
}
\]


\begin{voorbeeld}[Een verborgen tegenvoorbeeld]

Een leerling onderzoekt

\[
n^2>n
\]

voor natuurlijke getallen en controleert

\[
2^2>2,\qquad
3^2>3,\qquad
4^2>4,\qquad
5^2>5.
\]

De leerling besluit:

\begin{quote}
Voor elk natuurlijk getal $n$ geldt $n^2>n$.
\end{quote}

Die conclusie is te snel.

Als $0\in\mathbb{N}$ volgens de gebruikte conventie, dan geeft $n=0$

\[
0^2>0,
\]

wat onwaar is.

Ook voor $n=1$ krijgen we

\[
1^2>1,
\]

wat eveneens onwaar is.

Het voorbeeld toont meteen waarom zowel het \textbf{domein} als de
precieze formulering van een algemene uitspraak belangrijk zijn.
\end{voorbeeld}


\begin{oefening}[Zoek een tegenvoorbeeld]{\oefB \ESS}

Een leerling beweert:

\begin{center}
Elk oneven geheel getal is een priemgetal.
\end{center}

\begin{question}
Geef één tegenvoorbeeld.

\begin{oplossing}
Bijvoorbeeld $9$.

Het getal $9$ is oneven, maar

\[
9=3\cdot3,
\]

dus $9$ is geen priemgetal.
\end{oplossing}
\end{question}

\begin{question}
Waarom volstaat dit ene voorbeeld om de algemene bewering te
weerleggen?

\begin{oplossing}
De bewering zegt dat de eigenschap voor elk oneven geheel getal geldt.
Het getal $9$ behoort tot die verzameling maar bezit de beweerde
eigenschap niet. De algemene bewering kan daarom niet waar zijn.
\end{oplossing}
\end{question}

\end{oefening}


% ============================================================
\FVDsection{Logica in wetenschap en technologie}
% ============================================================

De begrippen uit dit hoofdstuk zijn niet alleen belangrijk binnen de
zuivere wiskunde.

Een computerprogramma, meettoestel of regelsysteem moet eveneens
beslissingen nemen op basis van voorwaarden die waar of onwaar zijn.


\begin{voorbeeld}[Een temperatuursensor]

Een koelsysteem moet inschakelen wanneer de temperatuur hoger wordt
dan $80\,{}^\circ\mathrm C$.

We kunnen de voorwaarde schrijven als

\[
P(T):\quad T>80,
\]

waarbij $T$ de temperatuur in graden Celsius voorstelt.

Bij

\[
T=72
\]

is $P(T)$ onwaar.

Bij

\[
T=93
\]

is $P(T)$ waar.

Een regelsysteem kan op basis van die waarheidswaarde een actie
uitvoeren.
\end{voorbeeld}


De instructie

\begin{quote}
Schakel de koeling in wanneer de temperatuur hoog is.
\end{quote}

is voor een computer onvoldoende precies.

Wat betekent immers \emph{hoog}?

De voorwaarde

\[
T>80\,{}^\circ\mathrm C
\]

legt dat ondubbelzinnig vast.


\begin{opmerking}
Dezelfde basisgedachte komt terug in onder andere

\begin{itemize}
  \item computerprogramma's;
  \item algoritmen;
  \item digitale schakelingen;
  \item automatische regelsystemen;
  \item wetenschappelijke modellen;
  \item databanken en zoekopdrachten.
\end{itemize}

Logica helpt ons niet alleen correct te rekenen.
Ze dwingt ons vooral om precies te formuleren
\textbf{wat we beweren}.
\end{opmerking}


\end{document}